\begin{flushright}
{\large ПРИЛОЖЕНИЕ 3}\\[0.5em]
Форма отзыва руководителя о практике
\end{flushright}

\vspace{1em}

\begin{center}
{\Large\bfseries ОТЗЫВ РУКОВОДИТЕЛЯ НАУЧНО-ИССЛЕДОВАТЕЛЬСКОЙ ПРАКТИКИ}
\end{center}

\noindent\textbf{Студент} \hfill \StudentName

\noindent\textbf{Университет} \hfill \UniversityName

\noindent\textbf{Факультет:} \FacultyName

\noindent\textbf{Направление:} \DirectionName

\noindent\textbf{Образовательная программа:} \ProgramName

\noindent\textbf{Место прохождения практики:} \PracticePlace

\noindent\textbf{Тема индивидуального задания:} \PracticeTheme

\vspace{1em}

\begin{center}
\textbf{ОЦЕНКА ДОСТИГНУТЫХ РЕЗУЛЬТАТОВ}
\end{center}

\small
\setlength{\tabcolsep}{3pt}
\begin{center}
\begin{tabularx}{\textwidth}{|C{1.1cm}|Y|C{2.2cm}|}
\hline
\textbf{№ п/п} & \textbf{Показатели$^{(*)}$} & \textbf{Оценка} \\
\hline
1. & Владение структурой исследовательского пайплайна и материалами дипломного проекта & зачтено \\
\hline
2. & Навыки подготовки данных, построения оконной выборки и контроля split/feature-контрактов & зачтено \\
\hline
3. & Анализ baseline-модели CatBoost и baseline-сравнений & зачтено \\
\hline
4. & Анализ промежуточных экспериментальных веток и корректность выводов о выборе финальной модели & зачтено \\
\hline
5. & Качество оформления итоговых материалов практики & зачтено \\
\hline
\multicolumn{2}{|c|}{\textbf{ИТОГОВАЯ ОЦЕНКА}} & зачтено \\
\hline
\end{tabularx}
\end{center}
\setlength{\tabcolsep}{6pt}
\normalsize

\vspace{1em}

\noindent\parbox{\textwidth}{$^{(*)}$ Перечисляются результаты оценивания, запланированные в рабочей программе практики.}

\vspace{1em}

\noindent\textbf{Отмеченные достоинства:}

В ходе научно-исследовательской практики студент последовательно проанализировал основные этапы проекта, посвященного разработке объяснимого прототипа рекомендательной системы выбора следующей культуры в севообороте. Работа выполнена с опорой на согласованный исследовательский пайплайн подготовки данных, формирования оконной выборки и построения baseline-модели CatBoost.

Студент корректно использовал результаты этапов очистки данных и построения оконного датасета, учел необходимость group-aware split по \texttt{CSBID} и аккуратно воспроизвел ключевые выводы по baseline-подходам и финальной модели \FinalModel{}. В отчете результаты проекта представлены с соблюдением исследовательской рамки: система описана как объяснимый рекомендательный прототип поддержки принятия решений.

\vspace{0.8em}

\noindent\textbf{Отмеченные недостатки:}

Существенных методических недостатков в представленных материалах не выявлено. В качестве возможного направления улучшения работы можно дополнительно рассмотреть другие классификационные модели и сравнить их с текущим baseline-подходом.

\vspace{0.8em}

\noindent\textbf{Заключение:}

Практика пройдена результативно. Поставленные исследовательские задачи выполнены, комплект отчетных материалов подготовлен, а выводы согласованы с содержанием дипломного проекта и текущими результатами репозитория.

\vspace{1.5em}

\noindent Руководитель практики

\vspace{0.5cm}

\noindent
\begin{tabularx}{\textwidth}{@{}L{5.2cm}C{3.2cm}Y@{}}
\SupervisorRole{} & \underline{\hspace{3cm}} & \SupervisorName
\end{tabularx}

\vspace{1.5cm}

\noindent «14» апреля 2026 г.
